% =====================================================================================================
% section2_expert_draft.tex — expert's draft of the formal pieces of Section 2 (2026-10-07)
%
% Three pieces, each a PASTE block for main.tex (line numbers refer to the Overleaf clone at 6cb4492):
%   PART A — Algorithm 1 (the master search as implemented) and the coupling paragraph.
%            Replaces lines 394 and 397–540; also gives replacement text for 2.2.1 (l. 550–577),
%            the lattice statement for 2.2.2–2.2.5 and the coupling paragraph 2.2.6 (l. 656–659).
%   PART B — new Subsection 2.2.7 "Recourse evaluation and certification".
%            Replaces 2.2.7 "Benders' Cuts" (l. 684–722) entirely.
%   PART C — Subsection 2.3 equation block: storage in the network models, cell-side throughput and the
%            degradation chain in the shared-ESS agent, the agent's objective, and the multi-scenario
%            commitment term. Replaces l. 735–906 (2.3.2 and the red note). The intro paragraphs of 2.3
%            (l. 724–733) are the author's; two sentences to change there are given at the top of Part C.
%
% Conventions: revised text is wrapped in \textcolor{blue}{...} as in the manuscript. Comments marked
% [CONFIRM] are points the Planner's equation-vs-code audit must settle before the text is final; they
% are LaTeX comments and do not print. Numbers are the frozen ones (tables v1, 590088fe).
% Nothing here is a result; results text belongs to Section 4.
% =====================================================================================================


% =====================================================================================================
% PART A — master search (Algorithm 1) and coupling
% =====================================================================================================

% ---- A.1 Replaces line 394 ("The two layers are coupled through the ESS investment variables. Benders'
% ---- decomposition iteratively updates ...") -----------------------------------------------------------

\textcolor{blue}{The two layers are coupled through the investment plan $\boldsymbol{x}$ only. The master layer
is a derivative-free search over the lattice of admissible plans: it uses the value $Q(\boldsymbol{x})$ returned
by the recourse and nothing else. No gradient, dual multiplier or cut is passed upward, because
$Q(\boldsymbol{x})$ is the value of a certified consensus point of a nonconvex problem (Subsection~2.2.7) and is
neither convex nor smooth in $\boldsymbol{x}$; the cut-based master of an earlier version of this work was
abandoned for that reason. To represent long-term storage ageing without sacrificing decomposition across
SOs, the operational layer includes a dedicated \textit{shared ESS agent} that tracks State-of-Health (SoH) and
the resulting evolution of available energy capacity. This preserves decomposition of the TSO and DSO
subproblems across representative years while maintaining intertemporal consistency of ESS degradation over
the planning horizon.}

% ---- A.2 Replaces the algorithm environment, lines 397–540 -------------------------------------------
% Notation: Delta^S = 0.25 MVA, Delta^E = 0.5 MWh (lattice units); n = number of lattice coordinates;
% Delta^p = poll size in lattice units; F(x) = I(x) + Q(x). Values of the parameters are given in Section 3.5.

\begin{algorithm}[htbp!]
    \caption{\textcolor{blue}{Derivative-free shared-ESS planning search (mesh adaptive direct search on the
    investment lattice) with certified ADMM recourse evaluations}}
    \label{alg:shared_ess_planning_mads}
    \DontPrintSemicolon

    \KwIn{Candidate nodes $E^S$; representative years $Y$, days $D$ and scenarios $\Omega_M$, $\Omega_O$ with
    probabilities; lattice units $\Delta^S$ (MVA) and $\Delta^E$ (MWh); duration bounds
    $\phi^{\text{Min}}, \phi^{\text{Max}}$; capacity limit $E^{\text{Max}}_{e,y}$; budget $c^{\text{Inv}}$;
    initial poll size $\Delta^p_0$; evaluation budget $N^{\max}$; certification parameters of
    Subsection~2.2.7; direction seed.}
    \KwOut{Incumbent plan $\boldsymbol{x}^{\star}$, its certified value $F(\boldsymbol{x}^{\star})$, the poll history
    and the evaluated neighbourhood with determinacy verdicts.}

    \tcp{Admissible set, enforced in closed form before any evaluation}
    $\mathcal{X} = \{\boldsymbol{x} = (S^{\text{Inv}}_{e,y}, E^{\text{Inv}}_{e,y})_{e \in E^S, y \in Y}:
    S^{\text{Inv}}_{e,y} \in \Delta^S \mathbb{Z}_{\ge 0},\ E^{\text{Inv}}_{e,y} \in \Delta^E \mathbb{Z}_{\ge 0},$
    \eqref{eq:master_duration}--\eqref{eq:master_budget}$\}$\;
    Scale: $\boldsymbol{z} = \mathrm{diag}(\Delta^S, \Delta^E, \ldots)^{-1}\boldsymbol{x}$, so that $\mathcal{X}$ is a
    subset of $\mathbb{Z}^{n}$ and the finest mesh is the lattice itself\;
    Cache $\mathcal{C} \gets \emptyset$ (one certified evaluation per canonical plan; a cache hit is never
    re-evaluated)\;

    \tcp{Initial search step: a designed set of lattice points (single-node ladders, combinations, staging)}
    Evaluate the designed set with Algorithm~\ref{alg:operational_planning_degradation} and the certification
    rule; add the records to $\mathcal{C}$\;
    $\boldsymbol{x}^{\text{inc}} \gets \arg\min_{\mathcal{C}} F$; $\Delta^p \gets \Delta^p_0$; $N \gets |\mathcal{C}|$\;

    \While{$N \le N^{\max}$}{
        \tcp{Poll step}
        Generate the $2n$ OrthoMADS directions $\pm \boldsymbol{v}_1, \ldots, \pm \boldsymbol{v}_n$ from the next
        Halton point of the recorded sequence\;
        Poll set $\mathcal{P} \gets \{\boldsymbol{z}^{\text{inc}} + \Delta^p \boldsymbol{v}\}$, each point rounded to the
        lattice and, if it violates \eqref{eq:master_duration}--\eqref{eq:master_budget}, snapped to the
        nearest admissible lattice point; drop duplicates and cached points\;
        \If{$|\mathcal{P}| < n + 1$}{
            Add to $\mathcal{P}$ the admissible lattice points within one unit step of $\boldsymbol{z}^{\text{inc}}$ in
            the $\infty$-norm (completion set), up to the stated cap\;
        }
        Evaluate every point of $\mathcal{P}$ (full poll, in batches of the available concurrent slots); an
        evaluation that does not certify, or whose storage agent is infeasible (SoH floor), is recorded with
        its cause and assigned $F = +\infty$ (extreme barrier); add the records to $\mathcal{C}$;
        $N \gets N + |\mathcal{P}|$\;
        \eIf{some $\boldsymbol{x} \in \mathcal{P}$ has $F(\boldsymbol{x}) < F(\boldsymbol{x}^{\text{inc}})$ by a
        determinate margin (Subsection~2.2.7)}{
            $\boldsymbol{x}^{\text{inc}} \gets$ the best such point; $\Delta^p \gets 2\Delta^p$\;
        }{
            \lIf{$\Delta^p = 1$}{\textbf{terminate}: $\boldsymbol{x}^{\text{inc}}$ is mesh-locally optimal on the lattice}
            $\Delta^p \gets \Delta^p / 2$\;
        }
    }
    \Return{$\boldsymbol{x}^{\star} = \boldsymbol{x}^{\text{inc}}$, $F(\boldsymbol{x}^{\star})$, the poll history, and the
    neighbourhood table of $\mathcal{C}$}\;
\end{algorithm}

% [CONFIRM — Planner, against p56a_oracle and the campaign launcher]
%  (a) the poll-size rule (double on a determinate success, halve on failure, floor 1) and the completion-set
%      cap (30) as run; (b) that "improvement" in the launcher is the determinacy rule of 2.2.7 and not a raw
%      comparison; (c) the Phase B variable form actually searched (single-cohort (S_e, E_e, y_e) per node,
%      with staging probed at the incumbent) — if so, add one sentence after the algorithm saying the
%      campaign searched that restriction of the lattice; (d) the number of evaluations and the seed, for 4.7.

% ---- A.3 Paragraph after the algorithm (new) ------------------------------------------------------------

\textcolor{blue}{Algorithm~\ref{alg:shared_ess_planning_mads} is a mesh adaptive direct search (MADS) on
granular variables~\cite{audet_dennis_2006,abramson_orthomads_2009,audet_granular_2019}. Every candidate is a
point of the investment lattice: converter power in multiples of $\Delta^S$, energy in multiples of
$\Delta^E$, within the duration, capacity and budget constraints of Subsection~2.2, which are enforced in closed
form before any evaluation. Each evaluation is one complete run of the operational recourse
(Algorithm~\ref{alg:operational_planning_degradation}, \ref{app:admm_updated_algorithm}), from a cold start and
under one frozen configuration, so that $Q(\boldsymbol{x})$ is a deterministic function of $\boldsymbol{x}$ and the
cache is exact. The search terminates when the poll at unit poll size fails, that is, when a positive spanning
set of lattice neighbours of the incumbent has been evaluated and none improves it by a determinate margin, or
when the evaluation budget is exhausted, in which case the incumbent is reported with the poll size reached.
The method therefore claims mesh-local optimality on the lattice at the stated resolution, and no global
optimality gap: the convergence theory of MADS applies to $F$ as a function on the lattice with the extreme
barrier and does not require continuity of $Q$, which is why a direct search was preferred to gradient- or
cut-based masters for this recourse.}

% [CONFIRM] bibliography keys: add to bibliography.bib —
%   audet_dennis_2006: C. Audet, J.E. Dennis Jr., "Mesh adaptive direct search algorithms for constrained
%     optimization", SIAM J. Optim. 17(1) 188–217, 2006.
%   abramson_orthomads_2009: M.A. Abramson, C. Audet, J.E. Dennis Jr., S. Le Digabel, "OrthoMADS: a
%     deterministic MADS instance with orthogonal directions", SIAM J. Optim. 20(2) 948–966, 2009.
%   audet_granular_2019: C. Audet, S. Le Digabel, C. Tribes, "The mesh adaptive direct search algorithm for
%     granular and discrete variables", SIAM J. Optim. 29(2) 1164–1189, 2019.

% ---- A.4 Replaces 2.2.1 "Objective Function", lines 550–577 ----------------------------------------------

\textcolor{blue}{The master problem minimizes the expected discounted investment cost plus the certified
expected operating cost of the plan:}

\textcolor{blue}{
    \begin{equation}
        \begin{aligned}
            \min_{\boldsymbol{x} \in \mathcal{X}} \; F(\boldsymbol{x}) &= I(\boldsymbol{x}) + Q(\boldsymbol{x}), \\
            I(\boldsymbol{x}) &=
            \sum_{e\in E^S} \sum_{y\in Y} \sum_{c\in\Omega_C}
            \frac{\omega_c}{(1+r)^{y-y_0}}
            \left(
                c^{\mathrm{Inv,S}}_{y,c} S^{\mathrm{Inv}}_{e,y}
                +
                c^{\mathrm{Inv,E}}_{y,c} E^{\mathrm{Inv}}_{e,y}
            \right)
        \end{aligned}
        \label{eq:master_objective}
    \end{equation}
}

\noindent
\textcolor{blue}{where $c^{\text{Inv,S}}_{y,c}$ and $c^{\text{Inv,E}}_{y,c}$ are the unit investment costs of
converter power rating and energy capacity under cost trajectory $c$, $S^{\text{Inv}}_{e,y}$ and
$E^{\text{Inv}}_{e,y}$ are the investment decisions, and $Q(\boldsymbol{x})$ is the expected discounted operating
cost of the plan returned by the recourse (Subsection~2.2.7). $I(\boldsymbol{x})$ is closed-form and cheap; every
evaluation cost is in $Q(\boldsymbol{x})$. The terminal (salvage) value of the installed assets is not part of
$F(\boldsymbol{x})$: operation does not internalize it, and it is reported beside every figure as a post-evaluated
credit computed from the terminal SoH and the remaining calendar life (Section~4).}

% [AUTHOR] 2.1, eq. (1) and the red salvage equation: make the sign a credit and state that it is reported,
% not optimized (the reviewers' R3.1 answer); the discount exponent should be y^End - y_0 with y^End the end
% of the last block.

% ---- A.5 Lattice statement for 2.2.2–2.2.5 (one sentence to add after the E/P constraint, l. 629; and the
% ---- labels the algorithm refers to) ----------------------------------------------------------------------
% Remove every "(l)" superscript and every "\forall l \in L^B" in 2.2.2–2.2.5: the plan is a point, not an
% iterate. Give the duration constraint the label \label{eq:master_duration} and the budget constraint the
% label \label{eq:master_budget}. Then add after line 629:

\textcolor{blue}{Investments are sized in commercial units: $S^{\text{Inv}}_{e,y}$ is a non-negative integer
multiple of $\Delta^S$ and $E^{\text{Inv}}_{e,y}$ of $\Delta^E$, and $S^{\text{Inv}}_{e,y} = 0$ implies
$E^{\text{Inv}}_{e,y} = 0$. The values used in the case study are $\Delta^S = 0.25$~MVA and
$\Delta^E = 0.5$~MWh with $\phi^{\text{Min}} = 2$~h and $\phi^{\text{Max}} = 4$~h (Section~3).}

% ---- A.6 Replaces 2.2.6 coupling paragraph, lines 656–659 ------------------------------------------------

\textcolor{blue}{For each candidate plan, the scenario-independent investment decisions $S^{\text{Inv}}_{e,y}$
and $E^{\text{Inv}}_{e,y}$ are passed to the operational recourse as fixed parameters; no probability-weighted
averaging of mutually exclusive physical ESS capacities is performed. The recourse returns the certified
value $Q(\boldsymbol{x})$ and its certification record; nothing else flows upward.}


% =====================================================================================================
% PART B — new Subsection 2.2.7
% =====================================================================================================
% Replaces \subsubsection{Benders' Cuts} and its three equations (l. 684–722).

\subsubsection{\textcolor{blue}{Recourse Evaluation and Certification}}
\label{subsubsec:certification}

\textcolor{blue}{The recourse value of a plan is the expected discounted operating cost of the coupled system
over the horizon,}

\textcolor{blue}{
    \begin{equation}
        Q(\boldsymbol{x}) =
        \sum_{y \in Y} \frac{Y_y}{(1+r)^{y-y_0}}
        \sum_{d \in D} D_d
        \sum_{m \in \Omega_M} \sum_{o \in \Omega_O} \omega_m \, \omega_o \,
        C^{\text{Op}}_{y,d,m,o}(\boldsymbol{x})
        \label{eq:recourse_value}
    \end{equation}
}

\noindent
\textcolor{blue}{where $C^{\text{Op}}_{y,d,m,o}$ is the daily operating cost of representative day $d$ of
representative year $y$ under market scenario $m$ and operation scenario $o$ at the consensus point
reached by the coordination procedure of \ref{app:admm_updated}: generation cost in the TN, activated
flexibility and curtailment in the DNs, and, in the multi-scenario instance, the commitment-deviation
charge of Subsection~2.3. Settlement flows between operators are transfers and are excluded from
$Q(\boldsymbol{x})$; they are reported separately. Each representative year stands for a block of $Y_y$
calendar years and is discounted at its representative year.}

% [CONFIRM] discount convention against shared_resources_planning.py (block weight Y_y at the representative
% year's discount factor — the "block-discounting approximation"); and the list of C^Op components against
% the penalty classification table (economic terms only).

\textcolor{blue}{\paragraph{Stopping and certification.} Each recourse evaluation runs the consensus ADMM of
\ref{app:admm_updated} until the primal and dual residuals of every consensus channel pass the absolute-plus-
relative test of Boyd et al.~\cite{boyd_admm_2011}, with $\varepsilon_{\text{abs}} = 10^{-5}$ and
$\varepsilon_{\text{rel}} = 10^{-4}$, in a cycle in which every local NLP solve is successful. Passing the
residual test does not imply that the objective has settled: on the reference evaluations of the case study the
objective moved by a further 15--21~k\euro{} after their residual-based certification. The residual pass at
cycle $k_0$ therefore only opens the certifying regime, which is then held fixed: acceleration off, the
tight interior-point tail on, and the penalty parameters $\rho$ frozen (\ref{app:admm_updated}). An
evaluation is certified at the first cycle $k^{\star} > k_0$ at which all of the following hold:}

\textcolor{blue}{
\begin{enumerate}
    \item the objective has shown at least three turning points since $k_0$, so that its period $\hat{P}$
    (measured between the first and third turning points) is measured rather than assumed;
    \item the successive half-swings of the objective are not growing; swings smaller than $\tau/10$ are
    treated as noise and are neither compared nor counted as turning points;
    \item the range of the objective over the last $W = \max\{20, \lceil 1.1\,\hat{P} \rceil\}$ cycles is at
    most $\tau$;
    \item the priced interface-consensus gap --- the value of the residual mismatch between the operators'
    interface schedules at the interface price --- is at most $\tau/2$ in absolute value;
    \item every cycle in that window was solved cleanly. A local solve is clean when it ends at the solver's
    optimal status, or at its acceptable status on a primary attempt with all four termination metrics within
    ten times the tight-tail tolerances.
\end{enumerate}
}

\textcolor{blue}{An objective that descends without oscillating is certified by a monotone branch when, over a
window of $2P_{\max}$ cycles with $P_{\max}$ the longest period measured on the instance, its steps are
decreasing, its range is at most $\tau$, and $|\text{last step}| \times 2P_{\max} \le \tau$, a linear bound on
the remaining descent that uses no fit. An evaluation that satisfies neither branch within its cycle cap is
reported as uncertified, with the objective's movement since the residual pass (its settling slack) and its
consensus gap.}

\textcolor{blue}{\paragraph{Resolution.} The threshold $\tau$ is set from the value resolution the paper claims:
$\tau = \delta_R V / 4$ with $\delta_R = 0.07$ and $V$ the storage value of the reference unit in force when the
rule was adopted ($V = 259{,}375.33$~\euro{}, hence $\tau = 4{,}539.07$~\euro{}; it was then kept frozen), so
that a ratio of two values --- four evaluations, each in error by at most $\tau$ --- is resolved to
$\delta_R$. The rule bounds the stopping error of each evaluation at about $\tau$, not well inside it: on the
case study ten of the certificates in use stop within 5\,\% of $\tau$, and evaluations continued past
certification moved by at most $0.9\,\tau$ (Section~4.7).}

\textcolor{blue}{\paragraph{Determinacy.} A difference between two certified evaluations is called determinate
when it exceeds $\max\{3\,b, 2\tau\}$, where $b$ is the larger of the two certification bands (the range of
the objective over the certifying window). A difference involving an uncertified evaluation is determinate
only if it exceeds three times the larger of that evaluation's consensus gap and settling slack. Every
planning statement in this paper is a pairwise comparison of evaluations with this threshold stated, and the
master search of Algorithm~\ref{alg:shared_ess_planning_mads} accepts a poll point as an improvement only by a
determinate margin; smaller differences are ties, reported as ``within resolution''. What is certified is
each evaluation, not the plan: the method claims no optimality certificate for the nonconvex recourse, and
the bands, gaps and slacks of every evaluation used in a reported comparison are tabulated in~\ref{app:results}.}

% [CONFIRM] (a) the settling slack definition as the scorer computes it (movement from the old certificate, or
% from the first residual pass where none exists) — Addendum 66(g); (b) "ten of the N certificates" — N for 4.7;
% (c) boyd_admm_2011 key exists in bibliography.bib (Boyd, Parikh, Chu, Peleato, Eckstein, Found. Trends ML
% 3(1), 2011; §3.3 for the stopping test, §3.4.1 for residual balancing).


% =====================================================================================================
% PART C — Subsection 2.3 equation block
% =====================================================================================================
% Two sentences of the author's 2.3 intro (l. 727 and l. 731) to change:
%   l. 727  "In this formulation, inter-year coupling arises exclusively from ESS degradation and is enforced
%           through the shared ESS subproblem."  → keep.
%   l. 731  add after "...across representative years.":
\textcolor{blue}{The storage schedule itself --- charging, discharging and reactive power in every hour --- is a
consensus variable between the operators' network models and the shared-ESS agent: the network models carry
the state-of-charge dynamics and the converter limits, and the agent carries the ageing chain that links the
representative years. The two views are reconciled by the ADMM consensus channel on the storage schedule
(\ref{app:admm_updated}).}

% ---- C.1 Replaces \subsubsection{Shared ESS Mathematical Formulation} (l. 735) and l. 737–790 -----------

\subsubsection{\textcolor{blue}{Installed and Available Capacity}}

\textcolor{blue}{The investment plan is a fixed parameter of the recourse. Each investment made at node $e$
in year $y^{\text{Inv}}$ forms a cohort that contributes its rated power and energy for its calendar
lifetime ($\forall e \in E^S$, $\forall y^\text{Inv} \in Y$, $\forall y \in [y^\text{Inv};~y^\text{Inv}+T^\text{Cal}_{e,y^\text{Inv}}]$):}

\textcolor{blue}{
    \begin{equation}
        S^{\text{Rated,Unit}}_{e,y^\text{Inv},y} = \hat{S}^{\text{Inv}}_{e,y^\text{Inv}}, \qquad
        E^{\text{Rated,Unit}}_{e,y^\text{Inv},y} = \hat{E}^{\text{Inv}}_{e,y^\text{Inv}}
        \label{eq:cohort_rated}
    \end{equation}
}

\textcolor{blue}{The available power of a cohort equals its rating, while its available energy is reduced by
its state of health at the start of the representative year ($\forall e$, $y^\text{Inv}$, $y$ as above):}

\textcolor{blue}{
    \begin{equation}
        S^{\text{Av,Unit}}_{e,y^\text{Inv},y} = S^{\text{Rated,Unit}}_{e,y^\text{Inv},y}, \qquad
        E^{\text{Av,Unit}}_{e,y^\text{Inv},y} = E^{\text{Rated,Unit}}_{e,y^\text{Inv},y}\, SoH_{e,y^\text{Inv},y-1}
        \label{eq:cohort_available}
    \end{equation}
}

% [CONFIRM] which SoH enters year y's available energy in the baseline: the start-of-block value SoH_{y-1}
% (written here; the mid-block evaluation point is the C3_midblock sensitivity arm) — check
% _configure_esso_cohort_state / get_updated_capacities.

\textcolor{blue}{and the capacities seen by the network models in year $y$ are the sums over the cohorts
alive in that year ($\forall e \in E^S$, $\forall y \in Y$):}

\textcolor{blue}{
    \begin{equation}
        S^{\text{Av}}_{e,y} = \sum_{y^\text{Inv}=\max(y-T^\text{Cal}_{e,y},\,y_0)}^{y} S^{\text{Av,Unit}}_{e,y^\text{Inv},y},
        \qquad
        E^{\text{Av}}_{e,y} = \sum_{y^\text{Inv}=\max(y-T^\text{Cal}_{e,y},\,y_0)}^{y} E^{\text{Av,Unit}}_{e,y^\text{Inv},y}
        \label{eq:available_capacity}
    \end{equation}
}

% ---- C.2 New: storage operation in the network models (answers R2.11, R2.12, R3.1, R3.3) ----------------

\subsubsection{\textcolor{blue}{Storage Operation in the Network Models}}

\textcolor{blue}{In every network model in which the shared ESS at node $e$ appears, its operation in
representative year $y$, day $d$ and hour $t$ is described by a charging power $P^{\text{Ch}}_{e,y,d,t} \ge 0$,
a discharging power $P^{\text{Dch}}_{e,y,d,t} \ge 0$, a reactive power $Q^{\text{E}}_{e,y,d,t}$ and a stored
energy $E^{\text{SoC}}_{e,y,d,t}$. The net active-power injection at the node is
$P^{\text{E}}_{e,y,d,t} = P^{\text{Dch}}_{e,y,d,t} - P^{\text{Ch}}_{e,y,d,t}$ (positive when discharging), and
the stored energy follows ($\forall e \in E^S$, $\forall y \in Y$, $\forall d \in D$, $\forall t \in T$):}

\textcolor{blue}{
    \begin{equation}
        E^{\text{SoC}}_{e,y,d,t} = E^{\text{SoC}}_{e,y,d,t-1}
        + \Delta t \left( \eta^{\text{Ch}}_e P^{\text{Ch}}_{e,y,d,t} - \frac{P^{\text{Dch}}_{e,y,d,t}}{\eta^{\text{Dch}}_e} \right)
        \label{eq:soc_recursion}
    \end{equation}
}

\textcolor{blue}{
    \begin{equation}
        SoC^{\text{Min}}_e \, E^{\text{Av}}_{e,y} \le E^{\text{SoC}}_{e,y,d,t} \le SoC^{\text{Max}}_e \, E^{\text{Av}}_{e,y}
        \label{eq:soc_limits}
    \end{equation}
}

\textcolor{blue}{
    \begin{equation}
        E^{\text{SoC}}_{e,y,d,t_0} = E^{\text{SoC}}_{e,y,d,|T|} = SoC^{0}_e \, E^{\text{Av}}_{e,y}
        \label{eq:soc_closure}
    \end{equation}
}

\textcolor{blue}{
    \begin{equation}
        P^{\text{Ch}}_{e,y,d,t} \le S^{\text{Av}}_{e,y}, \qquad
        P^{\text{Dch}}_{e,y,d,t} \le S^{\text{Av}}_{e,y}, \qquad
        \left(P^{\text{E}}_{e,y,d,t}\right)^2 + \left(Q^{\text{E}}_{e,y,d,t}\right)^2 \le \left(S^{\text{Av}}_{e,y}\right)^2
        \label{eq:converter_capability}
    \end{equation}
}

\textcolor{blue}{
    \begin{equation}
        P^{\text{Ch}}_{e,y,d,t} \, P^{\text{Dch}}_{e,y,d,t} = 0
        \label{eq:network_complementarity}
    \end{equation}
}

\noindent
\textcolor{blue}{where $\eta^{\text{Ch}}_e$ and $\eta^{\text{Dch}}_e$ are the charging and discharging
efficiencies, $SoC^{\text{Min}}_e$, $SoC^{\text{Max}}_e$ and $SoC^{0}_e$ the minimum, maximum and initial states
of charge as fractions of the available energy, and $\Delta t$ the hour length. Equation~\eqref{eq:soc_closure}
closes every representative day on its initial state, so no energy is carried between representative days and
cumulative discharge beyond the stored energy is excluded in every day and scenario by
\eqref{eq:soc_recursion}--\eqref{eq:soc_closure}. Reactive power is a converter quantity: it is limited with
the active power by the capability circle in \eqref{eq:converter_capability} and has no effect on the stored
energy or on ageing. The complementarity condition \eqref{eq:network_complementarity} is kept in the network
models in its relaxed, penalized form as in~\cite{simoes_2023}. In the multi-scenario instance the storage
schedule of a block is a single scenario-independent variable referenced by every operation scenario's
balance equations --- a day-ahead commitment --- so that the schedule the agent ages is the one every
scenario runs (Subsection~2.3.4).}

% [CONFIRM] against sess_soc_rule and the network ESS block: (a) the closure form (fixed initial value SoC^0,
% or free-but-equal first and last states); (b) whether the capability circle is on net P or on P^Ch and P^Dch
% separately; (c) the relaxed complementarity form and its penalty; (d) that the same block is instantiated in
% the TSO and in the DSO model of node e (consensus copies). Values of eta, SoC^Min/Max/0 go to Section 3.

% ---- C.3 Replaces l. 792–842 (degradation) ---------------------------------------------------------------

\subsubsection{\textcolor{blue}{Degradation Chain of the Shared-ESS Agent}}

\textcolor{blue}{The shared-ESS agent holds the consensus copy of the storage schedule and ages the cohorts
with it. Its schedule is allocated to the cohorts alive at node $e$ in proportion to their rated energy --- a
homogeneous-fleet approximation, in which all cohorts at a node cycle at the same depth ($\forall e$,
$y^\text{Inv}$, $y$, $d$, $t$):}

\textcolor{blue}{
    \begin{equation}
        P^{\text{Ch}}_{e,y^\text{Inv},y,d,t} = \frac{E^{\text{Rated,Unit}}_{e,y^\text{Inv},y}}{\sum_{y'} E^{\text{Rated,Unit}}_{e,y',y}}\, P^{\text{Ch}}_{e,y,d,t},
        \qquad
        P^{\text{Dch}}_{e,y^\text{Inv},y,d,t} = \frac{E^{\text{Rated,Unit}}_{e,y^\text{Inv},y}}{\sum_{y'} E^{\text{Rated,Unit}}_{e,y',y}}\, P^{\text{Dch}}_{e,y,d,t}
        \label{eq:cohort_allocation}
    \end{equation}
}

% [CONFIRM] H3 rule: the code allocates pnet_agg pro rata and lets the agent's own pch/pdch realize it; written
% here on P^Ch and P^Dch separately, which is equivalent once complementarity holds. Keep whichever the audit
% says the code does.

\textcolor{blue}{Ageing is driven by the energy that crosses the cells, not by the apparent power at the
converter: the average daily cell-side throughput of a cohort in representative year $y$ is}

\textcolor{blue}{
    \begin{equation}
        \bar{E}^{\text{Cell}}_{e,y^\text{Inv},y} =
        \sum_{d \in D} \frac{D_d}{365} \sum_{t \in T} \Delta t
        \left(
            \eta^{\text{Ch}}_e P^{\text{Ch}}_{e,y^\text{Inv},y,d,t}
            +
            \frac{P^{\text{Dch}}_{e,y^\text{Inv},y,d,t}}{\eta^{\text{Dch}}_e}
        \right)
        \label{eq:cell_throughput}
    \end{equation}
}

\textcolor{blue}{and the fractional cycling loss accumulated over the $Y_y$ calendar years of the block is}

\textcolor{blue}{
    \begin{equation}
        D_{e,y^\text{Inv},y} =
        \frac{365 \, Y_y \, \bar{E}^{\text{Cell}}_{e,y^\text{Inv},y}}
             {2 \, k_e \, E^{\text{Rated,Unit}}_{e,y^\text{Inv},y}}
        \label{eq:cycling_loss}
    \end{equation}
}

\noindent
\textcolor{blue}{where $k_e$ is the effective cycle life of the technology: $2 E^{\text{Rated,Unit}}$ is the
throughput of one full cycle, so $D$ counts equivalent full cycles against $k_e$. The state of health then
follows a continuous-compounding chain with a calendar retention factor $\phi^{\text{Cal}}_e$ per calendar
year ($\forall e$, $y^\text{Inv}$, $y$):}

\textcolor{blue}{
    \begin{equation}
        SoH_{e,y^\text{Inv},y} = SoH_{e,y^\text{Inv},y-1} \, \exp\!\left(-D_{e,y^\text{Inv},y}\right)
        \left(\phi^{\text{Cal}}_e\right)^{Y_y},
        \qquad SoH_{e,y^\text{Inv},y^\text{Inv}-1} = 1
        \label{eq:soh_chain}
    \end{equation}
}

\textcolor{blue}{
    \begin{equation}
        SoH_{e,y^\text{Inv},y} \ge SoH^{\text{Min}}_e
        \label{eq:soh_floor}
    \end{equation}
}

\textcolor{blue}{The constant $k_e$ is calibrated from a datasheet statement of $N^{\text{DS}}$ cycles at depth of
discharge $\delta^{\text{DS}}$ to a retention $R^{\text{DS}}$ of the rated capacity: with
\eqref{eq:cycling_loss}--\eqref{eq:soh_chain} and no calendar term, those cycles give
$SoH = \exp(-\delta^{\text{DS}} N^{\text{DS}} / k_e) = R^{\text{DS}}$, hence}

\textcolor{blue}{
    \begin{equation}
        k_e = \frac{\delta^{\text{DS}} N^{\text{DS}}}{-\ln R^{\text{DS}}}
        \label{eq:cycle_life_calibration}
    \end{equation}
}

\noindent
\textcolor{blue}{The calibrations used, the calendar retention and the end-of-life floor $SoH^{\text{Min}}_e$ are
given in Section~3.4. The floor is the only intertemporal trade-off enforced in operation: no explicit
cycling cost is charged, because a shared asset that has been paid for carries no operating cost of its own;
degradation is tracked and propagated, and its economic effect appears through the available energy of later
years and through the floor. Equations \eqref{eq:cycling_loss}--\eqref{eq:soh_floor} are linear in $D$ and
$SoH$ for a fixed schedule, and \eqref{eq:soh_chain} is the only nonlinear row of the agent.}

% [CONFIRM] (a) the row D*(2*cl_eff*E_rated) == 365*num_years*avg_ch_dch and the exp(-D)*phi^num_years chain
% as in shared_energy_storage_data.py; (b) that avg_ch_dch is cell-side (eta_ch*pch + pdch/eta_dch) and not
% terminal-side — Addendum/action plan record says cell-side; (c) num_years = block length of the
% representative year (equal 5-year blocks); (d) Section 3.4 values: C2 (10,000, 0.80, 0.80) → k = 35,851;
% C4 (8,000, 1.0, 0.70) → k = 22,430; phi_cal = 0.985; SoH_min = 0.70 (0.50 as a sensitivity row).

% ---- C.4 Replaces l. 844–900 (complementarity, net apparent power, objective) ---------------------------

\subsubsection{\textcolor{blue}{Objective of the Shared-ESS Agent}}

\textcolor{blue}{The agent has no economic objective. Within the coordination procedure it minimizes the
augmented-Lagrangian terms of the storage consensus channel (\ref{app:admm_updated_implementation}) plus a
small linear regularization of throughput,}

\textcolor{blue}{
    \begin{equation}
        \begin{aligned}
            \min \;
            & \varepsilon^{\text{E}} \sum_{e \in E^S} \sum_{y^\text{Inv} \in Y} \sum_{y \in Y} \sum_{d \in D} \sum_{t \in T}
            \left( P^{\text{Ch}}_{e,y^\text{Inv},y,d,t} + P^{\text{Dch}}_{e,y^\text{Inv},y,d,t} \right) \\
            & + \sum_{e \in E^S} \sum_{y \in Y} \sum_{d \in D} \sum_{t \in T}
            \left( \mathcal{L}^{E,P}_{e,y,d,t} + \mathcal{L}^{E,Q}_{e,y,d,t} \right)
        \end{aligned}
        \label{eq:esso_objective}
    \end{equation}
}

\noindent
\textcolor{blue}{subject to \eqref{eq:cohort_rated}--\eqref{eq:available_capacity},
\eqref{eq:cohort_allocation}--\eqref{eq:soh_floor}, $0 \le P^{\text{Ch}}_{e,y^\text{Inv},y,d,t},\,
P^{\text{Dch}}_{e,y^\text{Inv},y,d,t} \le S^{\text{Rated,Unit}}_{e,y^\text{Inv},y}$ and the converter capability
inequality of \eqref{eq:converter_capability} on the aggregate schedule. The regularization weight
$\varepsilon^{\text{E}}$ is small enough to be below the resolution of Subsection~2.2.7 and serves one purpose:
with a linear, strictly positive price on $P^{\text{Ch}} + P^{\text{Dch}}$, simultaneous charging and discharging
is never optimal, so complementarity in the agent follows from the problem structure and needs no
nonconvex constraint. The maximum residual $\min\{P^{\text{Ch}}, P^{\text{Dch}}\}$ is checked after every solve
and reported. The earlier investment-fixing slacks and the apparent-power aggregation are retired:
investments are parameters, and reactive power enters only through the capability circle.}

% [CONFIRM] (a) the value of EPS_ESSO_THROUGHPUT and its normalization (same units as the AL terms, which are
% normalized by the rating) — for Section 3.5; (b) whether the agent's capability circle is written on the
% aggregate or per cohort; (c) the post-solve detector threshold (1e-6 * s_max).

% ---- C.5 Replaces the red note (l. 905–906): the multi-scenario commitment term ------------------------

\subsubsection{\textcolor{blue}{Day-Ahead Commitment under Operation Scenarios}}
\label{subsubsec:commitment}

\textcolor{blue}{In the multi-scenario instance each operation scenario $o$ is instantiated inside every network
block with its own load and RES realization, while the interface schedule committed to the TSO is
scenario-independent by construction in the TSO model and may be deviated from in the DSO model. Each DSO
carries a day-ahead expectation of its interface flow and is charged for deviating from it. With
$P^{\text{I}}_{i,o,t}$ the interface active power of DSO $i$ in scenario $o$ and hour $t$ (likewise
$Q^{\text{I}}$), the expectation and the deviations are ($\forall i$, $o$, $t$):}

\textcolor{blue}{
    \begin{equation}
        \bar{P}^{\text{I}}_{i,t} = \sum_{o \in \Omega_O} \omega_o P^{\text{I}}_{i,o,t}, \qquad
        P^{\text{I}}_{i,o,t} - \bar{P}^{\text{I}}_{i,t} = d^{+}_{i,o,t} - d^{-}_{i,o,t}, \qquad d^{\pm}_{i,o,t} \ge 0
        \label{eq:commitment_deviation}
    \end{equation}
}

\textcolor{blue}{and the charge added to the DSO's objective, and hence to $C^{\text{Op}}$ in
\eqref{eq:recourse_value}, is}

\textcolor{blue}{
    \begin{equation}
        \sum_{o \in \Omega_O} \sum_{t \in T} \omega_o \, \alpha \, \bar{\pi}_t
        \left( d^{+}_{i,o,t} + d^{-}_{i,o,t} \right)
        \label{eq:commitment_charge}
    \end{equation}
}

\noindent
\textcolor{blue}{where $\bar{\pi}_t$ is the probability-weighted wholesale price of the hour over the market
scenarios and $\alpha$ an imbalance premium, so that a deviation is priced at a fraction $\alpha$ of the energy
price of the hour in which it occurs. The same term is written on the reactive interface flow with the same
price, as a stated assumption. The TSO operates on the committed schedule and carries no deviation term; the
interface voltage carries none either, being a physical state rather than a commitment. The storage schedule
is not deviated from: it is one scenario-free variable per hour in every block
(Subsection~2.3.2), so the committed schedule is the one that ages the cells. At a single operation scenario
\eqref{eq:commitment_deviation}--\eqref{eq:commitment_charge} vanish identically, and the single-scenario
instance is unaffected by them. The value of $\alpha$ is given in Section~3.5.}

% [CONFIRM] (a) alpha = 0.50 in the 3x3 campaign; (b) the expectation \bar P is a variable of the DSO block
% (its own day-ahead expectation), never the incoming consensus value; (c) the Q term uses the same alpha*pi;
% (d) label numbers: "Subsection 2.3.2" above refers to the "Storage Operation in the Network Models"
% subsubsection — fix to the final numbering.

% ---- C.6 Sentence to add to the degradation-attribution paragraph (l. 902), unchanged otherwise -----------
% Keep the blue paragraph at l. 902 as it is; it answers R1.2(v).

% =====================================================================================================
% NOMENCLATURE — symbols to add (Sets and Parameters / Variables tables); remove L^B, alpha^Down, alpha^(l),
% every "(l)" superscript, S^Ch/S^Dch (apparent charging/discharging), S^Net, P^Net/Q^Net, CL^Nom, delta_{e,y^Inv,y}.
%   Omega_M, omega_m            set of market scenarios and probability
%   Delta^S, Delta^E            lattice units of converter power (MVA) and energy (MWh)
%   Delta^p                     poll size of the master search, in lattice units
%   N^max                       evaluation budget of the master search
%   F(x), I(x), Q(x)            total, investment and certified operating cost of plan x
%   C^Op_{y,d,m,o}              daily operating cost of block (y,d) in scenario (m,o)
%   eps_abs, eps_rel            ADMM residual tolerances
%   k_0, k^*                    first residual pass; certification cycle
%   \hat P, P_max, W            measured period; longest period on the instance; certifying window
%   tau, delta_R                certification threshold; claimed value resolution
%   eta^Ch_e, eta^Dch_e         charging and discharging efficiencies
%   SoC^Min_e, SoC^Max_e, SoC^0_e   state-of-charge limits and initial state (fractions of available energy)
%   E^SoC_{e,y,d,t}             stored energy
%   P^Ch, P^Dch, Q^E, P^E       charging, discharging, reactive and net active power of the shared ESS
%   S^Av_{e,y}, E^Av_{e,y}      available power and energy at node e in year y
%   \bar E^Cell                 average daily cell-side throughput of a cohort
%   D_{e,y^Inv,y}               fractional cycling loss of a cohort over block y
%   k_e                         effective cycle life; N^DS, delta^DS, R^DS the datasheet calibration triple
%   phi^Cal_e                   calendar retention factor per year
%   SoH^Min_e                   end-of-life floor
%   eps^E                       throughput regularization weight of the agent
%   alpha, \bar pi_t            imbalance premium; expected hourly wholesale price
%   d^+, d^-                    positive and negative interface deviations from the day-ahead expectation
% =====================================================================================================


% =====================================================================================================
% PART D — §2.1, equations (1)–(2) and the salvage block (added 2026-10-07 after the author's agreement)
% Replaces, in the clone as pulled on 2026-10-07: from "\textcolor{blue}{The shared ESS planning problem
% can be generically written as:}" through the closing "}" of the blue "Here:" itemize (the red "and:" and
% the red salvage equation go). The paragraph "The absence of a scenario index..." that follows stays.
% =====================================================================================================

\textcolor{blue}{The shared ESS planning problem can be generically written as:}

\begin{equation}
    \textcolor{blue}{
        \min_{\boldsymbol{x}\in\mathcal{X}}
        \left\{
            I(\boldsymbol{x}) + Q(\boldsymbol{x})
        \right\}
    }
    \label{eq:two_stage_planning_problem}
\end{equation}

\noindent
\textcolor{blue}{where $I(\boldsymbol{x})$ is the expected discounted investment cost of the plan over the
investment-cost trajectories,}

\begin{equation}
    \textcolor{blue}{
        I(\boldsymbol{x}) = \sum_{c\in\Omega_C} \omega_c\, C_c^{\mathrm{Inv}}(\boldsymbol{x})
    }
    \label{eq:investment_cost_generic}
\end{equation}

\noindent
\textcolor{blue}{and $Q(\boldsymbol{x})$, the recourse function, is the expected operating cost of the plan under
coordinated TSO--DSO operation:}

\begin{subequations}
    \begin{align}
        \textcolor{blue}{Q(\boldsymbol{x})} &= \textcolor{blue}{\min_{\{\boldsymbol{u}_o\}_{o\in\Omega_O}} \sum_{o\in\Omega_O} \omega_o\, C_o^{\mathrm{Op}} \left(\boldsymbol{x}, \boldsymbol{u}_o \right)}\\
        \textcolor{blue}{\text{s.t.}} & \textcolor{blue}{\quad \boldsymbol{u}_o \in \mathcal{U}_o(\boldsymbol{x}), \quad \forall o\in\Omega_O}
    \end{align}
    \label{eq:operational_recourse_function}
\end{subequations}

\noindent
\textcolor{blue}{
    Here:
    \begin{itemize}
        \item $\mathcal{X}$ is the set of admissible plans (Subsection~\ref{subsec:master});
        \item $\Omega_C$ contains the investment-cost trajectories and $\omega_c$ is the probability of trajectory $c$;
        \item $\Omega_O$ contains the operating scenarios and $\omega_o$ is the probability of scenario $o$; in the
        multi-scenario instance a scenario is a pair of a market scenario $m \in \Omega_M$ and an operation scenario
        of load and RES generation, with probability $\omega_m \omega_o$, and the expansion of $Q(\boldsymbol{x})$
        over representative years, days and scenarios is given in Subsection~\ref{subsubsec:certification};
        \item $\boldsymbol{x}$ contains the scenario-independent investment decisions;
        \item $\boldsymbol{u}_o$ contains the operational variables for scenario $o$.
    \end{itemize}
}

\textcolor{blue}{The terminal value of the assets installed within the horizon is not part of
\eqref{eq:two_stage_planning_problem}: operation does not internalize it, and every planning result is
reported both gross and net of a salvage credit computed after the evaluation from the terminal state of
health and the remaining calendar life of each installation (Section~4).}

% ---- D.2 For Section 4, where the net columns are introduced (not in §2): the salvage credit ---------------
% Sign is a credit: net = gross − C^Salvage. [CONFIRM] against _get_terminal_salvage_value_*: (a) y^End is the
% end of the last block (2035 + Y_y) or the last representative year; (b) the SoH used is the terminal SoH of
% the last block; (c) whether the cohort's investment cost is the expected (probability-weighted) one.

\textcolor{blue}{
    \begin{equation}
        C^\text{Salvage}(\boldsymbol{x}) =
        \sum_{e \in E^S} \sum_{y^\text{Inv} \in Y}
        \frac{I_{e,y^\text{Inv}}}{(1+r)^{\,y^\text{End}-y_0}}
        \;\frac{T^\text{Rem}_{e,y^\text{Inv}}}{T^\text{Cal}_{e,y^\text{Inv}}}
        \; SoH_{e,y^\text{Inv},y^\text{End}}
        \label{eq:salvage_credit}
    \end{equation}
}

\textcolor{blue}{
    \begin{equation}
        T^\text{Rem}_{e,y^\text{Inv}} = \max\left\{0,\; y^\text{Inv} + T^\text{Cal}_{e,y^\text{Inv}} - y^\text{End}\right\}
        \label{eq:salvage_remaining_life}
    \end{equation}
}

\noindent
\textcolor{blue}{where $I_{e,y^\text{Inv}} = \sum_{c \in \Omega_C} \omega_c \big( c^{\text{Inv,S}}_{y^\text{Inv},c}
S^{\text{Inv}}_{e,y^\text{Inv}} + c^{\text{Inv,E}}_{y^\text{Inv},c} E^{\text{Inv}}_{e,y^\text{Inv}} \big)$ is the
expected undiscounted cost of the cohort, $y^\text{End}$ the end of the planning horizon, $T^\text{Rem}$ its
remaining calendar life at $y^\text{End}$, and $SoH_{e,y^\text{Inv},y^\text{End}}$ its terminal state of health.
The credit is reported, not optimized: the recourse does not see it, and a plan is never selected for it.}
